\subsection{纯 p-向量。格拉斯曼簇}

设 K 是一个域，E 是 K 上的一个向量空间。p-向量 \(z \in \bigwedge^{p}(\mathbf{E})\) 称为纯的（有时也称为可分解的），如果它是非零的并且存在 E 中的向量 \(x_{1}, \ldots, x_{p}\) 使得 \(z = x_{1} \wedge x_{2} \wedge \cdots \wedge x_{p}\) 。为此，必要且充分条件是：与 z 相伴的子空间 \(M_{z}\) （其维数总是 \(\geqslant p\) ，当 \(z \neq 0\) 时）恰好为 p 维（因为此时 \(\bigwedge^{p}(\mathbf{M}_{z})\) 的维数为 1）。特别地，每个非零标量、 \(\mathbf{E} = \bigwedge^{1}(\mathbf{E})\) 中的每个非零元素、以及当 E 的维数为 n 时 \(\bigwedge^{n}(\mathbf{E})\) 中的每个非零元素，都是纯的。

命题 15. 设 E 是一个 n 维向量空间，且设 e 为一个元素 \(\neq0\) （属于 \(\bigwedge^{n}(\mathbf{E})\) ，因而构成此向量空间的一组基）。设 \(\phi:\bigwedge(\mathbf{E})\to\bigwedge(\mathbf{E}^{*})\) 是与 e 相伴的向量空间同构（第 11 号，命题 12）。若 z 是 \(\bigwedge^{p}(\mathbf{E})\) 中的纯元素，则 \(\phi(z)\) 是 \(\bigwedge^{n-p}(\mathbf{E}^{*})\) 中的纯元素，且与 z 和 \(\phi(z)\) 相伴的子空间是正交的。

情形 \(p = 0\) 和 \(p = n\) 是平凡的。因此假设 \(1 \leqslant p \leqslant n - 1\) ，并设 \(z = x_1 \wedge \cdots \wedge x_p \neq 0\) 。于是存在一个基 \((e_i)_{1 \leqslant i \leqslant n}\) （ \(E\) 的），使得 \(e_i = x_i\) （对 \(1 \leqslant i \leqslant p\) ），且 \(e = e_1 \wedge e_2 \wedge \cdots \wedge e_n\) 。于是由第 11 号的公式 (78) 可得 \(\phi(z) = \pm e_{p+1}^* \wedge \cdots \wedge e_n^*\) ，由此得出该命题。

推论. 若 E 的维数为 n，则 E 的每个非零 \((n - 1)\) -向量都是纯的。

命题 16. 元素 \(z \neq 0\) （属于 \(\bigwedge^{p}(\mathbf{E})\) ）为纯的的必要充分条件是：对所有 \(u^{*} \in \bigwedge^{p-1}(\mathbf{E}^{*})\) ，

\[
(u ^ {*} \lrcorner z) \wedge z = 0.\tag{81}
\]

情形 \(p = 0\) 是平凡的，我们假设 \(p \geqslant 1\) 。若 \(z = x_1 \wedge \cdots \wedge x_p\) ，则公式 (68)（第 9 号）取 \(n = p - 1\) 表明 \(u^* \lrcorner z\) 是诸 \(x_i\) ( \(1 \leqslant i \leqslant p\) )的线性组合，由此得出 (81)。另一方面，若子空间 \(M_z\) （与 \(z\) 相伴）的维数 \(>p\) ，考虑此子空间的一组基 \((e_j)_{1 \leqslant j \leqslant n}\) （ \(n >p\) ）。由第 11 号命题 13 可知，每个 \(e_j\) 都具有 \(u^* \lrcorner z\) 的形式（对某个 \(u^* \in \bigwedge^{p-1}(E^*)\) ），因此关系式 (81) 蕴含 \(e_j \wedge z = 0\) （对 \(1 \leqslant j \leqslant n\) ）。由此可知，在表达式 \(z = \sum_{H} a_H e_H\) （其中 H 遍历 {[}1, n{]} 的含 \(p\) 个元素的子集的集合）中，所有系数 \(a_H\) 均为零，从而 \(z = 0\) ，与假设矛盾。

命题 16 的判据等价于当 \(u^*\) 遍历 \(\bigwedge^{p-1}(\mathbf{E}^*)\) 的一组基时写出条件 (81)。特别地，假设 \(\mathbf{E}\) 是有限维的，维数为 \(n\) ，并设 \((e_i)_{1\leq i\leq n}\) 是 \(\mathbf{E}\) 的一组基。则条件 (81) 等价于条件

\[
(82-(J, H)) \quad \langle e _ {\mathrm{J}} ^ {*}, (e _ {\mathrm{H}} ^ {*} \lrcorner z) \wedge z \rangle = 0
\]

对所有子集 J、H（ \([1, n]\) 的子集，满足 \(\text{Card}(J) = p + 1\) 和 \(\text{Card}(H) = p - 1\) ）成立。现在，若 I 和 \(I'\) 是 \([1, n]\) 的两个含 p 个元素的子集，则第 10 号的公式 (76) 以及 § 7 第 8 号的乘法表 (20) 表明

\[
\langle e _ {\mathrm{J}} ^ {*}, \langle e _ {\mathrm{H}} ^ {*} \lrcorner e _ {\mathrm{I}}) \wedge e _ {\mathrm{I} ^ {\prime}} \rangle = 0
\]

除非存在 \(i \in [1, n]\) ，使得 \(\mathbf{I} - \mathbf{H} = \{i\}\) 和 \(\mathbf{J} - \mathbf{I}' = \{i\}\) ，在这种情况下

\[
\langle e _ {\mathrm{J}} ^ {*}, (e _ {\mathrm{H}} ^ {*} \lrcorner e _ {\mathrm{I}}) \wedge e _ {\mathrm{I} ^ {\prime}} \rangle = (- 1) ^ {(p - 1) (p - 2) / 2} \varepsilon_ {i, \mathrm{J,H}}\tag{83}
\]

（其中 \(\varepsilon_{i,J,H} = \rho_{(i),H}\rho_{(i),I'}\) ；那么可以说，对 \(i\in J\cap \complement H\) , \(\varepsilon_{i,J,H}\) 等于 \(+1\) ，如果 \(J\) 中 \(< i\) 的元素个数与 \(H\) 中 \(< i\) 的元素个数有相同的奇偶性；否则等于 \(-1\) 。）

由此立即得出，若记 \(z = \sum_{\mathbf{I}} a_{\mathbf{I}} e_{\mathbf{I}}\) ，其中 \(\mathbf{I}\) 遍历 \([1, n]\) 的含 p 个元素的子集的集合，则关系 \((82-(J, H))\) 等价于关系

\[
(84 - (\mathrm{J}, \mathrm{H})) \quad \sum_ {i \in J \cup \complement_ {\mathrm{H}}} \varepsilon_ {i, J, \mathrm{H}} a _ {J - \{i \}} a _ {\mathrm{H} \cup \{i \}} = 0.
\]

关系 (84) 称为格拉斯曼关系；因此（当 J 遍历含 \(p + 1\) 个元素的子集的集合而 H 遍历含 \(p - 1\) 个元素的子集的集合（ \([1, n]\) 的）时），这些关系是元素 \(z \neq 0\) （属于 \(\bigwedge^{p}(\mathbf{E})\) ）为纯的必要充分条件。

注意，关系 (84) 并非独立的。例如，对 \(n = 4\) 和 \(p = 2\) ，格拉斯曼关系化简为单一关系

\[
a _ {1 2} a _ {3 4} - a _ {1 3} a _ {2 4} + a _ {1 4} a _ {2 3} = 0.\tag{85}
\]

设 \(D_p(E)\) 是 \(\bigwedge^p(E)\) 中由纯 \(p\) -向量组成的子集；显然 \(D_p(E)\) 关于 \(u\) 和 \(v\) 之间的等价关系：“存在 \(\lambda \in K^*\) 使得 \(v = \lambda u\)”是饱和的，且两个元素 \(u, v\) （ \(D_p(E)\) 的）在此关系下等价，当且仅当与它们相伴的子空间 \(M_u\) 和 \(M_v\) （ \(E\) 的）相同。因此我们由此得到一个从由 \(p\) 维向量子空间（ \(E\) 的）组成的集合到像 \(G_p(E)\) （即 \(D_p(E)\) 在射影空间 \(P(\bigwedge^p(E))\) （与 \(\bigwedge^p(E)\) 相伴）中的像）上的典范双射。子集 \(G_p(E)\) （ \(P(\bigwedge^p(E))\) 的）称为指标为 \(p\) 的向量空间 \(E\) 的格拉斯曼簇。当 \(E\) 有限维且 \((e_i)_{1 \leq i \leq n}\) 是 \(E\) 的一个基时，指标为 \(p\) 的格拉斯曼簇是 \(P(\bigwedge^p(E))\) 中使一个齐次坐标系 \((a_I)\) （相对于基 \((e_I)\) （ \(\bigwedge^p(E)\) 的基））满足格拉斯曼关系 (84) 的点的集合。

当 \(\mathbf{E} = \mathbf{K}^n\) 时，我们有时写 \(\mathbf{G}_{n,p}(\mathbf{K})\) 以代替 \(\mathbf{G}_p(\mathbf{K}^n)\) ，于是 \(\mathbf{G}_{n,1}(\mathbf{K}) = \mathbf{P}_{n-1}(\mathbf{K})\) 。映射 \(\mathbf{M} \mapsto \mathbf{M}^0\) 把每个 \(p\) 维子空间（ \(\mathbf{K}^n\) 的）对应到 \(\mathbf{E}^*\) 中的正交子空间（它借助于与 \(\mathbf{K}^n\) 的典范基对偶的基的选择而与 \(\mathbf{K}^n\) 等同），因而定义了 \(\mathbf{G}_{n,p}(\mathbf{K})\) 到 \(\mathbf{G}_{n,n-p}(\mathbf{K})\) 上的一个典范双射；命题 15 表明，这个双射是 \(\mathbf{G}_{n,p}(\mathbf{K})\) 上的这样一个限制：射影空间 \(\mathbf{P}(\bigwedge^p (\mathbf{K}^n))\) 到射影空间 \(\mathbf{P}(\bigwedge^{n-p}(\mathbf{K}^n))\) 的一个典范同构的限制。

